\addchaptertocentry{Удиртгал}

\titleformat{\section}{\Large\bfseries\filcenter}{}{0pt}{}
\section*{Удиртгал}
\titleformat{\section}{\Large\bfseries\raggedright}{\thesection}{1em}{}
\markboth{Удиртгал}{Удиртгал}

\textbf{Сэдвийн үндэслэл.} Банк бус санхүүгийн байгууллагад (ББСБ) зээлийн багц, хугацаа хэтрэлт, эргэн төлөлт,
итгэлцэл, ерөнхий дэвтрийн орлого зэрэг өдөр тутмын асуултад хариулахад шаардлагатай мэдээлэл өгөгдлийн санд
хадгалагддаг. Ийм асуултад
хариулахын тулд менежер, эрсдэлийн шинжээч, салбарын ажилтан SQL бичих чадвартай мэргэжилтэнд хүсэлт гаргаж,
хариугаа хүлээдэг. Байгалийн хэлний асуултыг өгөгдлийн сан дээр гүйцэтгэх SQL болгон хөрвүүлэх (Text2SQL)
системүүд энэ саадыг арилгах зорилготой бөгөөд том хэлний загвар (LLM) гарч ирснээр хурдацтай хөгжиж байна
\cite{shi_2025_survey_llm_text2sql}. Англи хэлний BIRD жишиг \cite{li_2023_bird} дээр нээлттэй эхийн 7 тэрбум
орчим параметртэй загварууд том хэмжээний синтетик өгөгдлөөр сургах \cite{li_2025_omnisql} болон бататгах сургалт
(RL) хийх \cite{yao_2025_arctic_text2sql_r1} замаар хаалттай том загваруудтай өрсөлдөх түвшинд хүрсэн.

Гэвч эдгээр үр дүн монгол хэлэнд шууд шилжихгүй. Олон хэлт судалгаанууд англи бус асуулт дээр загварын нарийвчлал
мэдэгдэхүйц буурдгийг харуулсан \cite{dou_2023_multispider, pham_2026_multispider2}. Монгол хэлний хувьд нэмэлт
гурван бэрхшээл бий. Нэгдүгээрт, залгамал үгзүйн улмаас нэг үг олон хэлбэртэй («зээл», «зээлийн», «зээлээс»).
Хоёрдугаарт, хэрэглэгчид кирилл болон латин галигийг ээлжлэн, заримдаа нэг асуултад холин бичдэг («zeeliin hvselt»,
«зээлийн хүсэлт», «mobile app-аар ирсэн хүсэлт»), галигийн хэв маяг нь ч нэг биш (ү → v/u, ө → u/o). Гуравдугаарт,
өгөгдлийн санд утга монголоор (кириллээр) хадгалагддаг тул асуултад галигаар эсвэл өөр хэлбэрээр бичигдсэн утгыг
хадгалсан утгатай нь яг тааруулах шаардлагатай. Олон хэлт жишгүүд энэ асуудлыг ихэвчлэн тойрдог: BIRDTurk асуулт ба
схемийг орчуулсан ч өгөгдлийн утгыг хөндөөгүй \cite{aktas_2026_birdturk}, IndicDB-д өгөгдөл англиар хадгалагдсан
бөгөөд буруу орчуулсан нэр шүүлтийг эвддэгийг тэмдэглэсэн \cite{dawar_2026_indicdb}. Монгол хэлэнд зориулсан
Text2SQL-ийн жишиг, үнэлгээний сан, туршилтын нотолгоо бараг байхгүй байгаа нь энэ ажлын үндэслэл болов.

Санхүүгийн салбарт нэмэлт шаардлага тавигдана: систем хувийн мэдээлэл задруулахгүй, өгөгдлийг өөрчлөхгүй, салбарын
эрхийн хязгаарыг сахих ёстой. Мөн хариулж чадахгүй асуултад буруу тоо гаргахаас илүү «мэдэхгүй» гэж хариулах нь
дээр \cite{lee_2024_trustsql}. Өгөгдлийг гадагш гаргахгүйн тулд загвар байгууллагын дотор (on-prem) ажиллах нь
чухал \cite{beskorovainyi_2026_onprem_frontier}.

\textbf{Өмнөх ажил.} Энэ ажлын өмнө цахим худалдааны өгөгдлийн сан дээр монгол/англи асуултыг SQLite, Oracle,
ClickHouse диалект руу хөрвүүлэх agent системийг (text2sql-thesis) туршсан \cite{text2sql_thesis_2026}. Тэрхүү
туршилтаар (dev, $n = 21$, хоёр ажиллуулалт) монгол асуулт англиасаа 24--29 нэгжээр хоцорсон боловч McNemar-ын
$p = 0.07$--$0.22$ тул статистикийн хувьд батлагдаагүй. Retrieval-ийн recall хоёр хэлэнд ижил байсан тул зөрүүг утгын
санамжийн нарийвчлал (хэсэгчлэн) ба SQL үүсгэлт (голлон) гэж задалсан. Түүвэр жижиг бөгөөд зэрэгцээ ажиллуулалт
детерминист бус байсан тул энэ дүгнэлтийг илүү том түүвэр дээр шалгах шаардлагатай байв. Энэ ажил нь тэрхүү
дүгнэлтээс үргэлжилж, (1) ББСБ-ийн үйл ажиллагааг загварчилсан схемтэй (синтетик өгөгдөлтэй) PostgreSQL сан,
(2) хариулах боломжгүй ба хоёрдмол асуулт бүхий том eval set, (3) fine-tuning ба agentic давхарга, (4) хоёр системийг
ижил нөхцөлд харьцуулах орчныг нэмж байна.

\textbf{Судалгааны зорилго.} Монгол хэлээр (кирилл, латин галиг) асуусан асуултыг ББСБ-ийн PostgreSQL өгөгдлийн сан
дээр гүйцэтгэх SQL болгон хөрвүүлэх, нээлттэй эхийн загварт суурилсан, байгууллагын дотор ажиллах системийг
судалгааны үр дүнд тулгуурлан бүтээж, аль арга монгол хэлэнд хамгийн их үр дүн өгөхийг хяналттай туршилтаар
тогтооход оршино.

\textbf{Судалгааны зорилтууд.}
\begin{enumerate}
  \item LLM-д суурилсан Text2SQL-ийн архитектур, RAG, fine-tuning, agent, бага нөөцтэй хэлний аргууд, үнэлгээний
        судалгааг тоймлох (Бүлэг~\ref{Chapter1});
  \item ББСБ-ийн өгөгдлийн сан, мэдлэгийн давхарга, хайлт, SQL хамгаалалт, загварын протокол, agentic урсгал,
        fine-tuning-ийн дамжлагыг зохиож хэрэгжүүлэх (Бүлэг~\ref{Chapter2});
  \item 500 асуулттай монгол eval set бүтээж, хариулах боломжгүй ба хоёрдмол асуултыг оруулах, gold SQL бүрийг
        гүйцэтгэлээр баталгаажуулах;
  \item Туршилтын төлөвлөгөө (A1--G6) боловсруулж, үнэлгээний дамжлагын зөв байдлыг шалгах, хоёр загварыг
        харьцуулах (Бүлэг~\ref{Chapter3});
  \item Өмнөх ажлын 150 кейсийг шинэ систем рүү шилжүүлж, хоёр системийг ижил загвар, өгөгдөл, шалгуураар
        харьцуулах.
\end{enumerate}

\textbf{Судалгааны асуултууд.}
\begin{description}
  \item[RQ1.] Нээлттэй эхийн Text2SQL загварыг монгол асуултад шууд ашиглахад англи хувилбараас хэр хоцордог вэ?
              Латин галиг нэмэлт хэр бууралт үүсгэх вэ? (A бүлэг)
  \item[RQ2.] Схемийн тайлбар, утгын хайлт, нэр томьёо, жишээний аль нь монгол хэлэнд хамгийн их нэмэр өгөх вэ? (B бүлэг)
  \item[RQ3.] Монгол өгөгдлөөр fine-tune хийх нь RAG-аас илүү үр дүнтэй юу, хоёуланг хослуулбал ямар вэ? (C, D бүлэг)
  \item[RQ4.] Засвар, санал хураалт, татгалзалт, agentic давталт нь найдвартай байдлыг (TrustSQL) хэр сайжруулах вэ? (E, G бүлэг)
  \item[RQ5.] Шинэ систем өмнөх agent системээс ижил нөхцөлд илүү үр дүнтэй юу? (thesis benchmark)
\end{description}

\textbf{Шинэлэг тал.} (1) Бидний мэдэхээр ББСБ-ийн схем дээрх анхны монгол Text2SQL eval set: кирилл, латин галиг, англи,
холимог хувилбартай, хариулах боломжгүй ба хоёрдмол асуулттай. (2) Монгол утгын хайлтыг кирилл ба галигийн сул
түлхүүрээр хийх, дагавар таслах дүрэм. (3) Бичгээс үл хамаарах reward (F4) зэрэг монгол хэлэнд чиглэсэн санаа, мөн
англи pivot аргыг (F5, F6) монгол Text2SQL-д тохируулсан хувилбар. (4) Гүйцэтгэлийн литерал шалгах tool-той agentic
давталт. (5) Хоёр системийг ижил загвар, өгөгдөл,
шалгуураар харьцуулах дахин давтагдах орчин.

\textbf{Ажлын бүтэц.} Бүлэг~\ref{Chapter1}-т Text2SQL-ийн онол ба судалгааны тоймыг, Бүлэг~\ref{Chapter2}-т системийн
зохиомж ба хэрэгжүүлэлтийг, Бүлэг~\ref{Chapter3}-т eval set, үнэлгээний арга, туршилтын төлөвлөгөө ба үр дүнг
тайлбарлана. Хавсралтад сервер дээр ажиллуулах заавар, eval set-ийн жишээ, тохиргооны жагсаалтыг оруулав. Системийг
«Монгол хэл дээрх Text2SQL: судалгаанд суурилсан бүтээх заавар» (цаашид «заавар») ба түүний туршилтын төлөвлөгөөний
(\tsql{docs/guide/}) дагуу бүтээв.
